
\documentclass{article}
\usepackage{colortbl}
\usepackage{makecell}
\usepackage{multirow}
\usepackage{supertabular}

\begin{document}

\newcounter{utterance}

\centering \large Interaction Transcript for game `dond', experiment `coop\_en', episode 8 with Qwen3.5{-}27B.
\vspace{24pt}

{ \footnotesize  \setcounter{utterance}{1}
\setlength{\tabcolsep}{0pt}
\begin{supertabular}{c@{$\;$}|p{.15\linewidth}@{}p{.15\linewidth}p{.15\linewidth}p{.15\linewidth}p{.15\linewidth}p{.15\linewidth}}
    \# & $\;$A & \multicolumn{4}{c}{Game Master} & $\;\:$B\\
    \hline

    \theutterance \stepcounter{utterance}  
    & & \multicolumn{4}{p{0.6\linewidth}}{
        \cellcolor[rgb]{0.9,0.9,0.9}{
            \makecell[{{p{\linewidth}}}]{
                \texttt{\tiny{[P1$\langle$GM]}}
                \texttt{You are playing a cooperative negotiation game in which you have to agree on how to divide a set of items among you and another player.} \\
\\ 
\texttt{Rules:} \\
\texttt{(a) You and the other player are given a set of items. Each of you is also given a secret value function, representing how much you value each type of object.} \\
\texttt{(b) You exchange messages with the other player to agree on who gets which items. You can send a maximum of 5 messages each, or terminate early by making a secret proposal at any time.} \\
\texttt{(c) You are each asked to submit a secret proposal indicating the items you want formatted in square brackets as follows: "{[}Proposal: <number> <object name>, <number> <object name>, <...>{]}"} \\
\texttt{(d) If your proposals are complementary, i.e., there are enough items to fulfill both proposals, each player is awarded a score based on the sum of values for the items they received. Otherwise, both of you get zero points.} \\
\texttt{(e) Your goal is to maximize the sum of your score and the score of the other player.} \\
\\ 
\texttt{Let us start.} \\
\\ 
\texttt{The set of available items is:} \\
\\ 
\texttt{2 potatoes, 1 glue, 1 berry, 3 carrots} \\
\\ 
\texttt{Your secret value function is:} \\
\\ 
\texttt{potato: 1, glue: 2, berry: 3, carrot: 1} \\
\\ 
\texttt{IMPORTANT: Your messages, unless it is the secret proposal, are directly transmitted to the other player, so do not include any response to the rules or text announcing your message. To make a secret proposal, use the indicated format. Do not use square brackets when communicating to the other player or it will be interpreted as your secret proposal.} \\
            }
        }
    }
    & \\ \\

    \theutterance \stepcounter{utterance}  
    & \multicolumn{4}{p{0.6\linewidth}}{
        \cellcolor[rgb]{0.9,0.9,0.9}{
            \makecell[{{p{\linewidth}}}]{
                \texttt{\tiny{[P1$\rangle$GM]}}
                \texttt{Hi! Let's figure out the best way to split these items to maximize our combined score. I have my own values, but I'm sure you have yours too.} \\
\\ 
\texttt{Here is what I have:} \\
\texttt{{-} 2 potatoes} \\
\texttt{{-} 1 glue} \\
\texttt{{-} 1 berry} \\
\texttt{{-} 3 carrots} \\
\\ 
\texttt{My values are: potato (1), glue (2), berry (3), carrot (1).} \\
\\ 
\texttt{Could you share your value function for these items? That way, we can see who values what the most and make sure the items go to the person who gets the most points from them.} \\
            }
        }
    }
    & & \\ \\

    \theutterance \stepcounter{utterance}  
    & & & \multicolumn{2}{p{0.3\linewidth}}{
        \cellcolor[rgb]{0.9,0.9,0.9}{
            \makecell[{{p{\linewidth}}}]{
                \texttt{\tiny{[GM$|$GM]}}
                \texttt{continue} \\
            }
        }
    }
    & & \\ \\

    \theutterance \stepcounter{utterance}  
    & & & \multicolumn{2}{p{0.3\linewidth}}{
        \cellcolor[rgb]{0.9,0.9,0.9}{
            \makecell[{{p{\linewidth}}}]{
                \texttt{\tiny{[GM$|$GM]}}
                \texttt{Hi! Let's figure out the best way to split these items to maximize our combined score. I have my own values, but I'm sure you have yours too.} \\
\\ 
\texttt{Here is what I have:} \\
\texttt{{-} 2 potatoes} \\
\texttt{{-} 1 glue} \\
\texttt{{-} 1 berry} \\
\texttt{{-} 3 carrots} \\
\\ 
\texttt{My values are: potato (1), glue (2), berry (3), carrot (1).} \\
\\ 
\texttt{Could you share your value function for these items? That way, we can see who values what the most and make sure the items go to the person who gets the most points from them.} \\
            }
        }
    }
    & & \\ \\

    \theutterance \stepcounter{utterance}  
    & & \multicolumn{4}{p{0.6\linewidth}}{
        \cellcolor[rgb]{0.9,0.9,0.9}{
            \makecell[{{p{\linewidth}}}]{
                \texttt{\tiny{[GM$\rangle$P2]}}
                \texttt{You are playing a cooperative negotiation game in which you have to agree on how to divide a set of items among you and another player.} \\
\\ 
\texttt{Rules:} \\
\texttt{(a) You and the other player are given a set of items. Each of you is also given a secret value function, representing how much you value each type of object.} \\
\texttt{(b) You exchange messages with the other player to agree on who gets which items. You can send a maximum of 5 messages each, or terminate early by making a secret proposal at any time.} \\
\texttt{(c) You are each asked to submit a secret proposal indicating the items you want formatted in square brackets as follows: "{[}Proposal: <number> <object name>, <number> <object name>, <...>{]}"} \\
\texttt{(d) If your proposals are complementary, i.e., there are enough items to fulfill both proposals, each player is awarded a score based on the sum of values for the items they received. Otherwise, both of you get zero points.} \\
\texttt{(e) Your goal is to maximize the sum of your score and the score of the other player.} \\
\\ 
\texttt{Let us start.} \\
\\ 
\texttt{The set of available items is:} \\
\\ 
\texttt{2 potatoes, 1 glue, 1 berry, 3 carrots} \\
\\ 
\texttt{Your secret value function is:} \\
\\ 
\texttt{potato: 1, glue: 0, berry: 2, carrot: 2} \\
\\ 
\texttt{IMPORTANT: Your messages, unless it is the secret proposal, are directly transmitted to the other player, so do not include any response to the rules or text announcing your message. To make a secret proposal, use the indicated format. Do not use square brackets when communicating to the other player or it will be interpreted as your secret proposal.} \\
\\ 
\\ 
\texttt{Hi! Let's figure out the best way to split these items to maximize our combined score. I have my own values, but I'm sure you have yours too.} \\
\\ 
\texttt{Here is what I have:} \\
\texttt{{-} 2 potatoes} \\
\texttt{{-} 1 glue} \\
\texttt{{-} 1 berry} \\
\texttt{{-} 3 carrots} \\
\\ 
\texttt{My values are: potato (1), glue (2), berry (3), carrot (1).} \\
\\ 
\texttt{Could you share your value function for these items? That way, we can see who values what the most and make sure the items go to the person who gets the most points from them.} \\
            }
        }
    }
    & \\ \\

    \theutterance \stepcounter{utterance}  
    & & & \multicolumn{4}{p{0.6\linewidth}}{
        \cellcolor[rgb]{0.9,0.9,0.9}{
            \makecell[{{p{\linewidth}}}]{
                \texttt{\tiny{[GM$\langle$P2]}}
                \texttt{Hi! Thanks for sharing your values. Here is my value function:} \\
\\ 
\texttt{{-} potato: 1} \\
\texttt{{-} glue: 0} \\
\texttt{{-} berry: 2} \\
\texttt{{-} carrot: 2} \\
\\ 
\texttt{Looking at our values together:} \\
\texttt{{-} **Potatoes**: We both value them at 1. It doesn't matter who gets them for the total score, but we need to split the 2 available.} \\
\texttt{{-} **Glue**: You value it at 2, I value it at 0. You should definitely take the glue.} \\
\texttt{{-} **Berry**: You value it at 3, I value it at 2. You should take the berry.} \\
\texttt{{-} **Carrots**: You value them at 1, I value them at 2. I should take the carrots.} \\
\\ 
\texttt{To maximize our combined score, the optimal split seems to be:} \\
\texttt{{-} **You**: 1 glue, 1 berry, and maybe some potatoes.} \\
\texttt{{-} **Me**: 3 carrots, and the remaining potatoes.} \\
\\ 
\texttt{Since there are 2 potatoes and we both value them equally, we can split them 1 each, or you can take both if you prefer, as long as we don't overlap on the high{-}value items.} \\
\\ 
\texttt{How does this sound? If you agree, I can make a secret proposal for:} \\
\texttt{{-} Me: 3 carrots, 1 potato} \\
\texttt{{-} You: 1 glue, 1 berry, 1 potato} \\
\\ 
\texttt{Or if you prefer to take both potatoes, just let me know!} \\
            }
        }
    }
     \\ \\

    \theutterance \stepcounter{utterance}  
    & & & \multicolumn{2}{p{0.3\linewidth}}{
        \cellcolor[rgb]{0.9,0.9,0.9}{
            \makecell[{{p{\linewidth}}}]{
                \texttt{\tiny{[GM$|$GM]}}
                \texttt{continue} \\
            }
        }
    }
    & & \\ \\

    \theutterance \stepcounter{utterance}  
    & & & \multicolumn{2}{p{0.3\linewidth}}{
        \cellcolor[rgb]{0.9,0.9,0.9}{
            \makecell[{{p{\linewidth}}}]{
                \texttt{\tiny{[GM$|$GM]}}
                \texttt{Hi! Thanks for sharing your values. Here is my value function:} \\
\\ 
\texttt{{-} potato: 1} \\
\texttt{{-} glue: 0} \\
\texttt{{-} berry: 2} \\
\texttt{{-} carrot: 2} \\
\\ 
\texttt{Looking at our values together:} \\
\texttt{{-} **Potatoes**: We both value them at 1. It doesn't matter who gets them for the total score, but we need to split the 2 available.} \\
\texttt{{-} **Glue**: You value it at 2, I value it at 0. You should definitely take the glue.} \\
\texttt{{-} **Berry**: You value it at 3, I value it at 2. You should take the berry.} \\
\texttt{{-} **Carrots**: You value them at 1, I value them at 2. I should take the carrots.} \\
\\ 
\texttt{To maximize our combined score, the optimal split seems to be:} \\
\texttt{{-} **You**: 1 glue, 1 berry, and maybe some potatoes.} \\
\texttt{{-} **Me**: 3 carrots, and the remaining potatoes.} \\
\\ 
\texttt{Since there are 2 potatoes and we both value them equally, we can split them 1 each, or you can take both if you prefer, as long as we don't overlap on the high{-}value items.} \\
\\ 
\texttt{How does this sound? If you agree, I can make a secret proposal for:} \\
\texttt{{-} Me: 3 carrots, 1 potato} \\
\texttt{{-} You: 1 glue, 1 berry, 1 potato} \\
\\ 
\texttt{Or if you prefer to take both potatoes, just let me know!} \\
            }
        }
    }
    & & \\ \\

    \theutterance \stepcounter{utterance}  
    & & \multicolumn{4}{p{0.6\linewidth}}{
        \cellcolor[rgb]{0.9,0.9,0.9}{
            \makecell[{{p{\linewidth}}}]{
                \texttt{\tiny{[P1$\langle$GM]}}
                \texttt{Hi! Thanks for sharing your values. Here is my value function:} \\
\\ 
\texttt{{-} potato: 1} \\
\texttt{{-} glue: 0} \\
\texttt{{-} berry: 2} \\
\texttt{{-} carrot: 2} \\
\\ 
\texttt{Looking at our values together:} \\
\texttt{{-} **Potatoes**: We both value them at 1. It doesn't matter who gets them for the total score, but we need to split the 2 available.} \\
\texttt{{-} **Glue**: You value it at 2, I value it at 0. You should definitely take the glue.} \\
\texttt{{-} **Berry**: You value it at 3, I value it at 2. You should take the berry.} \\
\texttt{{-} **Carrots**: You value them at 1, I value them at 2. I should take the carrots.} \\
\\ 
\texttt{To maximize our combined score, the optimal split seems to be:} \\
\texttt{{-} **You**: 1 glue, 1 berry, and maybe some potatoes.} \\
\texttt{{-} **Me**: 3 carrots, and the remaining potatoes.} \\
\\ 
\texttt{Since there are 2 potatoes and we both value them equally, we can split them 1 each, or you can take both if you prefer, as long as we don't overlap on the high{-}value items.} \\
\\ 
\texttt{How does this sound? If you agree, I can make a secret proposal for:} \\
\texttt{{-} Me: 3 carrots, 1 potato} \\
\texttt{{-} You: 1 glue, 1 berry, 1 potato} \\
\\ 
\texttt{Or if you prefer to take both potatoes, just let me know!} \\
            }
        }
    }
    & \\ \\

    \theutterance \stepcounter{utterance}  
    & \multicolumn{4}{p{0.6\linewidth}}{
        \cellcolor[rgb]{0.9,0.9,0.9}{
            \makecell[{{p{\linewidth}}}]{
                \texttt{\tiny{[P1$\rangle$GM]}}
                \texttt{That sounds like a perfect plan! Your analysis is spot on.} \\
\\ 
\texttt{{-} **Glue**: You correctly identified I should take it (my value 2 vs your 0).} \\
\texttt{{-} **Berry**: I should take it (my value 3 vs your 2).} \\
\texttt{{-} **Carrots**: You should take all 3 (your value 2 vs my 1).} \\
\texttt{{-} **Potatoes**: Since we both value them at 1, splitting them 1 each is a fair and efficient way to use them up.} \\
\\ 
\texttt{So, the final agreement is:} \\
\texttt{{-} **Me**: 1 glue, 1 berry, 1 potato.} \\
\texttt{{-} **You**: 3 carrots, 1 potato.} \\
\\ 
\texttt{This maximizes our total score. I will now submit my secret proposal to lock this in.} \\
\\ 
\texttt{{[}Proposal: 1 glue, 1 berry, 1 potato{]}} \\
            }
        }
    }
    & & \\ \\

    \theutterance \stepcounter{utterance}  
    & & & \multicolumn{2}{p{0.3\linewidth}}{
        \cellcolor[rgb]{0.9,0.9,0.9}{
            \makecell[{{p{\linewidth}}}]{
                \texttt{\tiny{[GM$|$GM]}}
                \texttt{proposal} \\
            }
        }
    }
    & & \\ \\

    \theutterance \stepcounter{utterance}  
    & & & \multicolumn{2}{p{0.3\linewidth}}{
        \cellcolor[rgb]{0.9,0.9,0.9}{
            \makecell[{{p{\linewidth}}}]{
                \texttt{\tiny{[GM$|$GM]}}
                [1, 1, 1, 0]
            }
        }
    }
    & & \\ \\

    \theutterance \stepcounter{utterance}  
    & & \multicolumn{4}{p{0.6\linewidth}}{
        \cellcolor[rgb]{0.9,0.9,0.9}{
            \makecell[{{p{\linewidth}}}]{
                \texttt{\tiny{[GM$\rangle$P2]}}
                \texttt{The other player has submitted a secret proposal. It is now your turn to submit a proposal of your own.} \\
            }
        }
    }
    & \\ \\

    \theutterance \stepcounter{utterance}  
    & & & \multicolumn{4}{p{0.6\linewidth}}{
        \cellcolor[rgb]{0.9,0.9,0.9}{
            \makecell[{{p{\linewidth}}}]{
                \texttt{\tiny{[GM$\langle$P2]}}
                \texttt{{[}Proposal: 1 potato, 3 carrots{]}} \\
            }
        }
    }
     \\ \\

    \theutterance \stepcounter{utterance}  
    & & & \multicolumn{2}{p{0.3\linewidth}}{
        \cellcolor[rgb]{0.9,0.9,0.9}{
            \makecell[{{p{\linewidth}}}]{
                \texttt{\tiny{[GM$|$GM]}}
                \texttt{proposal} \\
            }
        }
    }
    & & \\ \\

    \theutterance \stepcounter{utterance}  
    & & & \multicolumn{2}{p{0.3\linewidth}}{
        \cellcolor[rgb]{0.9,0.9,0.9}{
            \makecell[{{p{\linewidth}}}]{
                \texttt{\tiny{[GM$|$GM]}}
                [1, 0, 0, 3]
            }
        }
    }
    & & \\ \\

    \theutterance \stepcounter{utterance}  
    & & & \multicolumn{2}{p{0.3\linewidth}}{
        \cellcolor[rgb]{0.9,0.9,0.9}{
            \makecell[{{p{\linewidth}}}]{
                \texttt{\tiny{[GM$|$GM]}}
                [[1, 1, 1, 0], [1, 0, 0, 3]]
            }
        }
    }
    & & \\ \\

\end{supertabular}
}

\end{document}
